% Осень26!B20; нумерация внутри столбца без пустых строк.
% Источник: лекция 3, §6 (строка 168)
\definitionticket{17}{Формулировка теоремы о лебеговском продолжении меры}
\label{definition:17}

\begin{theorem}
Пусть $\mathcal P$ --- полукольцо подмножеств $X$, а $\mu_0$ --- $\sigma$-конечная мера на $\mathcal P$.
Тогда существуют $\sigma$-алгебра $\mathcal M$ и мера $\mu$ на ней со следующими свойствами:
\begin{enumerate}
  \item $\mathcal P\subset\mathcal M$ и $\mu|_{\mathcal P}=\mu_0$.
  \item Мера $\mu$ полна.
  \item Если $\mu'$ --- полная мера на $\sigma$-алгебре $\mathcal M'\supset\mathcal P$ и $\mu'|_{\mathcal P}=\mu_0$, то $\mathcal M\subset\mathcal M'$ и $\mu'|_{\mathcal M}=\mu$.
  \item Если $\mathcal P\subset\mathcal P'\subset\mathcal M$, $\mathcal P'$ --- полукольцо и $\mu'$ --- мера на $\mathcal P'$, продолжающая $\mu_0$, то $\mu' = \mu|_{\mathcal P'}$.
  \item Для любого $A\in\mathcal M$
  \[
    \mu(A)=\inf\left\{\sum_{i=1}^{\infty}\mu_0(P_i):
    \ A\subset\bigcup_{i=1}^{\infty}P_i,\ P_i\in\mathcal P\right\}.
  \]
\end{enumerate}
\end{theorem}
